\section{Molekulalar massasi. Modda miqdori}


\begin{enumerate}

\item 16 g kislorod moddasining miqdorini (mol) aniqlang. $M = 32 \cdot 10^{-3}$ kg/mol.
\begin{enumerate}[label=\Alph*)]
\item 0,5
\item 1
\item 2
\item 16
\end{enumerate}

\item 450 mol kislorodning massasini aniqlang (kg).
\begin{enumerate}[label=\Alph*)]
\item 4,5
\item 7,2
\item 9
\item 14,4
\end{enumerate}

\item Modda miqdorining birligi nima?
\begin{enumerate}[label=\Alph*)]
\item kg
\item N
\item mol
\item kg/mol
\end{enumerate}

\item Modda miqdorlari teng bo‘lgan qalay va qo‘rg‘oshindan qilingan jismlarning massalarini va hajmlarini taqqoslang. $M_{Sn} = 116,7$ g/mol, $\rho_{Sn} = 7300$ kg/$m^3$, $M_{Pb} = 207,2$ g/mol, $\rho_{Pb} = 11300$ kg/$m^3$.
\begin{enumerate}[label=\Alph*)]
\item Qo'rg'oshin jismning massasi 1,7 marta hajmi 1,1 marta katta
\item Qalay jismning massasi 1,7 marta hajmi 1,1 marta katta
\item Qo'rg'oshin jismning massasi 1,7 marta hajmi 1,1 marta kichik
\item Qalay jismning massasi 1,7 marta hajmi 1,1 marta kichik
\end{enumerate}

\item 1 g karbonat angidrid ($CO_2$) da qancha molekula bor?
\begin{enumerate}[label=\Alph*)]
\item $1,4 \cdot 10^{22}$
\item $1,4 \cdot 10^{19}$
\item $2,8 \cdot 10^{22}$
\item $1,2 \cdot 10^{22}$
\end{enumerate}

\item Diametri 1 mm li suv tomchisida nechta molekula bor? $N_A = 6,02 \cdot 10^{23}$ $mol^{-1}$.
\begin{enumerate}[label=\Alph*)]
\item $1,75 \cdot 10^{18}$
\item $2 \cdot 10^{17}$
\item $1,75 \cdot 10^{29}$
\item $1,75 \cdot 10^{19}$
\end{enumerate}

\item Molekulalar soni $3 \cdot 10^{23}$ bo‘lgan azotning massasini hisoblang (g)? $M = 28$ g/mol.
\begin{enumerate}[label=\Alph*)]
\item 14
\item 28
\item 56
\item 84
\end{enumerate}

\item Massasi 135 g bo‘lgan alyuminiy buyumda qancha atom borligini toping. $M_{Al} = 27$ g/mol.
\begin{enumerate}[label=\Alph*)]
\item $2 \cdot 10^{24}$
\item $6 \cdot 10^{24}$
\item $1,5 \cdot 10^{24}$
\item $3 \cdot 10^{24}$
\end{enumerate}

\item 36 $sm^3$ hajmdagi suvda nechta molekula bor? $M = 18$ g/mol; $N_A = 6 \cdot 10^{23}$ $mol^{-1}$.
\begin{enumerate}[label=\Alph*)]
\item $3,6 \cdot 10^{24}$
\item $2,4 \cdot 10^{24}$
\item $1,2 \cdot 10^{24}$
\item $6 \cdot 10^{24}$
\end{enumerate}

\item Massasi 3 g bo‘lgan suv tomchisida nechta elektron bor? $N_A = 6 \cdot 10^{23}$ $mol^{-1}$.
\begin{enumerate}[label=\Alph*)]
\item $10^{26}$
\item $10^{25}$
\item $10^{24}$
\item $10^{23}$
\end{enumerate}

\item Berk idishdagi gaz molekulalari soni 68 marta kamaytirildi. Gazning massasi qanday o‘zgaradi?
\begin{enumerate}[label=\Alph*)]
\item 68 marta kamayadi
\item 68 marta ortadi
\item 34 marta kamayadi
\item 34 marta ortadi
\end{enumerate}

\item Massasi 20 g bo‘lgan sof oltin uzukda nechta oltin atomi bo‘ladi? Oltinning molyar massasi 197 g/mol ga teng. $N_A = 6 \cdot 10^{23}$ 1/mol.
\begin{enumerate}[label=\Alph*)]
\item $6,1 \cdot 10^{22}$
\item $5,1 \cdot 10^{22}$
\item $6,7 \cdot 10^{22}$
\item $3,1 \cdot 10^{23}$
\end{enumerate}

\item 3 l suvdagi molekulalar sonini aniqlang. Avogadro soni $6 \cdot 10^{23}$ $mol^{-1}$.
\begin{enumerate}[label=\Alph*)]
\item $1 \cdot 10^{25}$
\item $3,3 \cdot 10^{25}$
\item $3,3 \cdot 10^{26}$
\item $1 \cdot 10^{26}$
\end{enumerate}

\item Moddaning molyar massasi XBS da qanday birlikka ega?
\begin{enumerate}[label=\Alph*)]
\item $kg \cdot m^{-3} \cdot mol$
\item $kg \cdot mol$
\item $kg \cdot m^{-3}$
\item $kg \cdot mol^{-1}$
\end{enumerate}

\item Molekulyar kinetik nazariyaning asosiy qoidalari qaysi javobda nisbatan to‘liq bayon qilingan?
\begin{enumerate}[label=\Alph*)]
\item Har qanday jism molekulalardan tuzilgan, ular betartib harakatda bo'ladi, ular orasida o'zaro ta'sir kuchlari mavjud.
\item Modda mayda zarralardan tashkil topgan va ular orasida o'zaro ta'sir kuchlari mavjud.
\item Modda mayda bo'laklardan tuziladi va ular fazoni uzluksiz to'ldiradi.
\item Modda elementar zarralardan tuziladi va ular bir-biriga tinimsiz aylanib turadi.
\end{enumerate}

\item Oltinning nisbiy atom massasi 197 g/mol ga teng. Massasi 3,94 g bo‘lgan oltin uzukda nechta oltin atomi bor? $N_A = 6 \cdot 10^{23}$ $mol^{-1}$.
\begin{enumerate}[label=\Alph*)]
\item $2,4 \cdot 10^{22}$
\item $6 \cdot 10^{21}$
\item $1,1 \cdot 10^{21}$
\item $1,2 \cdot 10^{22}$
\end{enumerate}

\item $m = 49$ g bo‘lgan sulfat kislotalarda ($H_2SO_4$) nechta molekula bor? Atomning nisbiy massalari: $A_r(H)=1, A_r(S)=32$ va $A_r(O)=16, N_A = 6 \cdot 10^{23}$ $mol^{-1}$.
\begin{enumerate}[label=\Alph*)]
\item $10^{23}$
\item $3 \cdot 10^{23}$
\item $4 \cdot 10^{23}$
\item $5 \cdot 10^{23}$
\end{enumerate}

\item $m = 96$ g metanda ($CH_4$) qancha molekula borligini aniqlang. $N_A = 6 \cdot 10^{23}$ $mol^{-1}$.
\begin{enumerate}[label=\Alph*)]
\item $10^{23}$
\item $3,6 \cdot 10^{24}$
\item $3,6 \cdot 10^{23}$
\item $5 \cdot 10^{23}$
\end{enumerate}

\item Zichligi 4,4 kg/$m^3$ bo‘lgan karbonat angidrid gazining 1 $m^3$ hajmida qancha molekula bor? $M = 44 \cdot 10^{-3}$ kg/mol; $N_A = 6 \cdot 10^{23}$ $mol^{-1}$.
\begin{enumerate}[label=\Alph*)]
\item $1 \cdot 10^{25}$
\item $1,1 \cdot 10^{24}$
\item $4 \cdot 10^{25}$
\item $6 \cdot 10^{25}$
\end{enumerate}

\item 1 g to‘la dissotsiyalangan azotda molekulalar soni qancha bo‘lishini aniqlang. $M_{Az} = 28$ g/mol.
\begin{enumerate}[label=\Alph*)]
\item $26 \cdot 10^{22}$
\item $43 \cdot 10^{22}$
\item $2,6 \cdot 10^{22}$
\item $4,3 \cdot 10^{22}$
\end{enumerate}

\item Avogadro soni deb qanday fizik kattalikka aytiladi?
\begin{enumerate}[label=\Alph*)]
\item 12 g ugleroddagi atomlar soniga
\item 1 mol moddadagi zarralar soniga
\item 32 g kisloroddagi molekulalar soniga
\item 2 g vodoroddagi molekulalar soniga
\end{enumerate}

\item Molyar massa deb nimaga aytiladi?
\begin{enumerate}[label=\Alph*)]
\item $0^\circ C$ da 1 $m^3$ hajmdagi modda massasiga
\item 1 mol moddaning massasiga
\item Uglerod atomi massasining 1/12 qismiga nisbatiga
\item $N_A = 6 \cdot 10^{23}$ ta zarrachadan tashkil topgan modda massasiga
\end{enumerate}

\item Uranning nisbiy atom massasi 238 m.a.b. ga teng. Bitta uran atomining massasini aniqlang (g). 1 m.a.b. = $1,66 \cdot 10^{-27}$ kg.
\begin{enumerate}[label=\Alph*)]
\item $3,95 \cdot 10^{-23}$
\item $3,95 \cdot 10^{-22}$
\item $238 \cdot 10^{-24}$
\item $3,95 \cdot 10^{-26}$
\end{enumerate}

\end{enumerate}



\begin{table}[!ht]
\centering
\begin{tabular}{|c|c|c|c|c|c|c|c|c|c|}
\hline
№ & J. & № & J. & № & J. & № & J. & № & J. \\ \hline
1 & A & 6 & D & 11 & A & 16 & D & 21 & B \\ \hline
2 & D & 7 & A & 12 & A & 17 & B & 22 & B \\ \hline
3 & C & 8 & D & 13 & D & 18 & B & 23 & B \\ \hline
4 & A & 9 & C & 14 & D & 19 & B & & \\ \hline
5 & A & 10 & C & 15 & A & 20 & B & & \\ \hline
\end{tabular}
\end{table}